\documentclass[11pt,a4paper]{article}

% --- Variables du TP
\newcommand{\TPModule}{Datawarehouse}
\newcommand{\TPNumero}{2}
\newcommand{\TPKeywords}{datawarehouse, datamart, star schema, dbt, KPI}
\newcommand{\TPSujet}{Modélisation décisionnelle et couche GOLD (CORRECTION)}
\newcommand{\TPTheme}{Clinique}
\usepackage[utf8]{inputenc}
\usepackage[french]{babel}
\usepackage[T1]{fontenc}
\usepackage{amsmath}
\usepackage{amsfonts}
\usepackage{xcolor}
\usepackage{amssymb}
\usepackage{graphicx}
\usepackage[left=2cm,right=2cm,top=2cm,bottom=2cm]{geometry}
\usepackage{fancyhdr}
\usepackage[bottom]{footmisc}
\usepackage{tikz}
\usepackage{fourier}
\usepackage{pst-text}
\usepackage{multicol}
\usepackage{comment}
\usepackage{array}
\usepackage{enumitem}
\setlist[itemize]{label=\textbullet} % pour ne plus utiliser les bulletpoints après \item
\setlength{\parindent}{0pt} % pour supprimer les \noindent
\usepackage{amssymb} % pour les symboles
\usepackage{tabularx}
\usepackage{tcolorbox}
\usepackage{fontawesome}
\usepackage{listings}

% --- Réglages listings (inchangé)
\lstset{
  basicstyle=\ttfamily\small,
  backgroundcolor=\color{gray!10},
  frame=single,
  rulecolor=\color{gray!40},
  breaklines=true,
  columns=fullflexible,
  showstringspaces=false,
  aboveskip=0.5em,
  belowskip=0.5em,
  keywordstyle=\color{blue},
  commentstyle=\color{gray},
}
\lstdefinelanguage{markdown}{
  sensitive=true,
  morecomment=[l]{<!--},
  morecomment=[s]{<!--}{-->},
  morekeywords={},
}

% --- Réglages listings (inchangé)
\lstset{
  basicstyle=\ttfamily\small,
  backgroundcolor=\color{gray!10},
  frame=single,
  rulecolor=\color{gray!40},
  breaklines=true,
  columns=fullflexible,
  showstringspaces=false,
  aboveskip=0.5em,
  belowskip=0.5em,
  keywordstyle=\color{blue},
  commentstyle=\color{gray},
}
\lstdefinelanguage{markdown}{
  sensitive=true,
  morecomment=[l]{<!--},
  morecomment=[s]{<!--}{-->},
  morekeywords={},
}

% --- Fancyhdr construit depuis les variables
\pagestyle{fancy}
\setlength{\headheight}{26pt}
\renewcommand{\baselinestretch}{1.3}

\fancyhead[L]{\textbf{EPSI\\\TPModule}}
\fancyhead[R]{\textbf{TP n°\TPNumero\\\TPSujet}}
\renewcommand{\arraystretch}{1}

% --- Hyperref + metadata construits depuis les variables
\usepackage[
  colorlinks=true,
  linkcolor=blue!70!black,
  urlcolor=blue!70!black,
  citecolor=blue!70!black
]{hyperref}

\hypersetup{
  pdftitle={TP n°\TPNumero\space - \TPSujet},
  pdfauthor={Thibault CLEMENT},
  pdfsubject={TP n°\TPNumero\space - \TPSujet},
  pdfkeywords={\TPKeywords},
}

% --- Composants réutilisables
\newcommand{\TPFrontPage}{%
  ~
  \begin{center}
    \framebox{\parbox[c]{11cm}{%
      \begin{center}
        \Large{\textbf{TP n°\TPNumero~:~\TPSujet\\\large{\TPTheme}}}
      \end{center}
    }}
  \end{center}
  \vspace{0.5cm}
}

% --- Tes macros custom existantes
\newcommand{\file}[1]{~\texttt{\textcolor{blue!70!black}{#1}}~}

\newtcolorbox{reflectionbox}[1][]{%
  colback=blue!5,
  colframe=blue!60,
  title={\faLightbulbO~~Réflexions},
  fonttitle=\bfseries,
  left=1mm,
  right=1mm,
  top=1mm,
  bottom=1mm,
  #1
}


\newcommand{\checkbox}{\(\square\)~~}
\graphicspath{{images/}}
\setlength{\columnseprule}{1pt}
\def\columnseprulecolor{\color{black}}



\begin{document}
\TPFrontPage

% ---------------------------------------------------------------------------------------------

\section*{Contexte}

Lors du TP1, vous avez :

\begin{itemize}
  \item exploré les données sources,
  \item mis en place une couche \textbf{BRONZE} historisée,
  \item construit une couche \textbf{SILVER} propre et testée via dbt.
\end{itemize}

Vous devez maintenant concevoir la \textbf{couche GOLD},
c’est-à-dire un \textbf{datamart décisionnel} permettant de répondre
aux questions stratégiques du directeur.

\bigskip

\section*{Rappel : Questions métier du directeur}

L’objectif final est de pouvoir répondre aux questions suivantes :

\begin{enumerate}
  \item Quels services sont le plus souvent en tension ?
  \item Quel est le taux moyen d’occupation par service ?
  \item Combien de patients sont refusés chaque semaine et pourquoi ?
  \item La satisfaction diminue-t-elle en période de forte activité ?
  \item Existe-t-il un lien entre présence du personnel et tension ?
  \item Quelle est la durée moyenne de séjour par service ?
  \item Quels services semblent les plus rentables ?
  \item Peut-on anticiper les semaines critiques ?
\end{enumerate}

\textbf{Objectif du TP2 :}
construire le modèle décisionnel permettant de produire ces indicateurs.

% ---------------------------------------------------------------------------------------------

\section*{Architecture technique}

Base de données : \file{dwh/clinic\_dwh.db}

Séparation des couches par préfixe :

\begin{itemize}
  \item \file{bronze\_*}
  \item \file{silver\_*}
  \item \file{gold\_*} (à construire)
\end{itemize}

Les modèles GOLD devront être implémentés via \textbf{dbt}
dans un dossier \file{models/gold/}.

% ---------------------------------------------------------------------------------------------

\section{Étape 1 : Définition du fait principal et du grain}

\subsection*{Objectif}

Définir précisément ce que l’on souhaite analyser.

\subsection*{Travail demandé}

\begin{enumerate}
  \item Identifier le \textbf{fait métier principal}.
    \begin{tcolorbox}[title=\textbf{Fait métier principal}, colback=gray!5, colframe=black, boxrule=0.8pt, arc=2mm]
    L’activité hebdomadaire d’un service de la clinique.
    \end{tcolorbox}
  \item Définir son \textbf{grain exact}.
    \begin{tcolorbox}[title=\textbf{Grain choisi}, colback=gray!5, colframe=black, boxrule=0.8pt, arc=2mm]
      \begin{center}
      1 ligne = 1 service × 1 semaine (yearweek)
      \end{center}
    \end{tcolorbox}
  \item Justifier ce choix.
    \begin{tcolorbox}[title=\textbf{Justification}, colback=gray!5, colframe=black, boxrule=0.8pt, arc=2mm]
    \begin{itemize}
    \item Les questions du directeur portent sur la \textbf{performance hebdomadaire des services}.
    \item Le niveau semaine permet d’analyser la tension, la saisonnalité et les pics d’activité.
    \item Ce grain évite les problèmes de duplication liés aux données transactionnelles.
    \item Il constitue un bon compromis entre granularité fine et lisibilité décisionnelle.
    \end{itemize}
    \end{tcolorbox}
  \item Lister les mesures associées.
    \begin{tcolorbox}[title=\textbf{Mesures associées}, colback=gray!5, colframe=black, boxrule=0.8pt, arc=2mm]
    \begin{itemize}
    \item requests\_count
    \item admitted\_count
    \item refused\_count
    \item refusal\_rate
    \item patients\_present\_count (proxy basé sur arrival\_yearweek)
    \item avg\_los
    \item avg\_satisfaction
    \item staff\_present\_count
    \end{itemize}
    \end{tcolorbox}
\end{enumerate}

\textbf{Indication :}
un grain pertinent est souvent :
\begin{center}
\textit{1 ligne = 1 service × 1 semaine}
\end{center}


% ---------------------------------------------------------------------------------------------

\section{Étape 2 : Construction des indicateurs hebdomadaires}

Aucune table agrégée hebdomadaire n’est fournie.
Vous devez la construire à partir de la couche \textbf{silver}.

\subsection*{Objectif}

Créer une table décisionnelle agrégée par \textbf{service} et \textbf{yearweek}.

\subsection*{Indicateurs attendus (minimum)}

\begin{itemize}
  \item Nombre de patients présents.
  \item Nombre de demandes.
  \item Nombre de patients admis.
  \item Nombre de patients refusés.
  \item Taux de refus.
  \item Durée moyenne de séjour.
  \item Satisfaction moyenne.
  \item Nombre de membres du personnel présents.
\end{itemize}

\subsection*{Travail demandé}

\begin{enumerate}
  \item Identifier les tables silver nécessaires.
  \item Gérer les jointures et éventuelles duplications.
  \item Créer un modèle dbt : \file{gold\_fact\_service\_weekly}.
  \item Documenter les calculs effectués.
\end{enumerate}

\subsection*{Points d’attention}

\begin{itemize}
  \item Attention au grain lors des jointures.
  \item Vérifier la cohérence : 
  \file{patients\_refused = requests - admitted}.
  \item Ne pas utiliser d’information future pour calculer une métrique passée.
\end{itemize}

\begin{tcolorbox}[title=\textbf{Gestion du grain et des jointures :}, colback=gray!5, colframe=black, boxrule=0.8pt, arc=2mm]

\textbf{Tables silver utilisées :}
\begin{itemize}
\item silver\_admissions\_requests
\item silver\_patients
\item silver\_staff\_schedule
\item silver\_staff
\end{itemize}

\textbf{Principe fondamental :}
Les agrégations sont réalisées \textbf{avant les jointures}.  
Chaque source est d’abord agrégée au grain (service × yearweek).

\textbf{Pourquoi ?}
\begin{itemize}
\item Éviter les duplications (multiplication des lignes).
\item Garantir la cohérence des KPI.
\item Respecter strictement le grain défini.
\end{itemize}

\textbf{Contrôle métier :}
On vérifie systématiquement :
\[
refused\_count = requests\_count - admitted\_count
\]

\end{tcolorbox}

% ---------------------------------------------------------------------------------------------

\section{Étape 3 : Construction des dimensions}

\subsection*{Objectif}

Mettre en place un schéma en étoile.

\subsection*{Dimensions minimales attendues}

\begin{itemize}
  \item \file{gold\_dim\_service}
  \item \file{gold\_dim\_date} (ou week)
  \item \file{gold\_dim\_staff} (si pertinent)
\end{itemize}

\subsection*{Travail demandé}

\begin{enumerate}
  \item Définir les attributs de chaque dimension.
  \item Identifier les clés primaires.
  \item Relier les dimensions à la table de faits.
  \item Justifier les cardinalités.
\end{enumerate}

\begin{tcolorbox}[title=\textbf{Schéma en étoile proposé :}, colback=gray!5, colframe=black, boxrule=0.8pt, arc=2mm]

\textbf{Table de faits :}
\begin{itemize}
\item gold\_fact\_service\_weekly
\end{itemize}

\textbf{Dimensions :}
\begin{itemize}
\item gold\_dim\_service (1 → n)
\item gold\_dim\_week (1 → n)
\item gold\_dim\_staff (optionnelle pour analyses RH)
\end{itemize}

\textbf{Cardinalités :}
\begin{itemize}
\item 1 service → plusieurs lignes de faits
\item 1 semaine → plusieurs lignes de faits
\end{itemize}

Ce modèle respecte les principes du \textbf{star schema} :
\begin{itemize}
\item table centrale contenant les mesures,
\item dimensions descriptives autour,
\item simplicité pour les outils BI.
\end{itemize}

\end{tcolorbox}

\begin{tcolorbox}[title=\textbf{Pourquoi ne pas avoir créé de dimension patient ?}, colback=gray!5, colframe=black, boxrule=0.8pt, arc=2mm]

Il est volontaire que le modèle GOLD ne contienne pas de dimension \texttt{patient}.

\textbf{Justification :}

\begin{itemize}
\item Le grain de la table de faits est :  
\begin{center}
1 ligne = 1 service × 1 semaine
\end{center}

\item À ce niveau d’agrégation, le patient n’est plus une entité d’analyse individuelle, mais une contribution à un indicateur agrégé (comptage, moyenne, etc.).

\item Introduire une dimension patient impliquerait :
  \begin{itemize}
  \item soit de descendre le grain au niveau \textit{patient × semaine},
  \item soit de créer un autre datamart orienté analyse individuelle.
  \end{itemize}

\item Les questions du directeur portent sur la \textbf{performance des services}, pas sur le suivi individuel des patients.
\end{itemize}

\textbf{Conclusion :}  
Dans ce datamart, le patient est une \textbf{mesure agrégée}, pas une dimension d’analyse.

\end{tcolorbox}

\begin{tcolorbox}[colback=blue!5, colframe=blue!60, title=\textbf{Commentaire}]
Un modèle décisionnel doit être conçu en fonction des questions métier.
On ne modélise pas toutes les entités disponibles,
mais uniquement celles pertinentes au niveau d’analyse choisi.
\end{tcolorbox}

% ---------------------------------------------------------------------------------------------

\section{Étape 4 : Implémentation GOLD avec dbt}

\subsection*{Consignes}

\begin{enumerate}
  \item Implémenter les modèles GOLD dans \file{models/gold/}.
  \item Matérialisation en tables.
  \item Documenter les modèles.
\end{enumerate}

\subsection*{Tests obligatoires}

\begin{itemize}
  \item \file{unique} + \file{not\_null} sur les PK.
  \item \file{relationships} entre faits et dimensions.
  \item Test métier simple sur un KPI.
\end{itemize}

\subsection*{Commandes attendues}

\begin{lstlisting}
dbt run  --select gold
dbt test --select gold
\end{lstlisting}

\begin{tcolorbox}[title=\textbf{Stratégie de tests dbt :}, colback=gray!5, colframe=black, boxrule=0.8pt, arc=2mm]

\textbf{Tests techniques :}
\begin{itemize}
\item unique + not\_null sur les clés primaires.
\item relationships entre faits et dimensions.
\end{itemize}

\textbf{Test métier :}
Vérification de la cohérence d’un KPI :
\[
refused = requests - admitted
\]

Ces tests garantissent :
\begin{itemize}
\item l’intégrité référentielle,
\item la cohérence des indicateurs,
\item la fiabilité du datamart.
\end{itemize}

\end{tcolorbox}

% ---------------------------------------------------------------------------------------------

\section{Étape 5 : Historisation et volumétrie}

\subsection*{Réflexion}

\begin{enumerate}
  \item Estimer la volumétrie annuelle.
    \begin{tcolorbox}[title=\textbf{Volumétrie estimée :}, colback=gray!5, colframe=black, boxrule=0.8pt, arc=2mm]

    La volumétrie d’un datamart en étoile est principalement portée par la
    \textbf{table de faits}.

    Grain choisi :
    \[
    1\ service \times 1\ semaine
    \]

    Avec 4 services :
    \[
    4 \times 52 = 208\ lignes/an
    \]

    Même sur 10 ans :
    \[
    208 \times 10 = 2080\ lignes
    \]

    La couche GOLD reste donc extrêmement peu volumineuse.

    Les dimensions (service, week, staff) évoluent très peu
    et n’ont pas d’impact significatif sur la volumétrie globale.
    \end{tcolorbox}

    \begin{tcolorbox}[colback=blue!5, colframe=blue!60, title=\textbf{Commentaire enseignant}]
    La volumétrie dépend directement du grain choisi.
    Si le grain avait été :
    1 ligne = 1 patient × 1 jour,
    la volumétrie aurait été plusieurs ordres de grandeur plus élevée.

    Le choix du grain impacte donc :
    \begin{itemize}
    \item la taille des tables,
    \item les performances,
    \item la complexité du modèle.
    \end{itemize}
    \end{tcolorbox}

  \item Identifier une dimension susceptible d’évoluer.
    \begin{tcolorbox}[title=\textbf{Dimension susceptible d’évoluer :}, colback=gray!5, colframe=black, boxrule=0.8pt, arc=2mm]
    La dimension staff (changement de rôle, de service).
    \end{tcolorbox}
  \item Proposer une stratégie SCD (Type 1 ou Type 2). Justifier le choix.
    \begin{tcolorbox}[title=\textbf{Stratégie SCD proposée et justification :}, colback=gray!5, colframe=black, boxrule=0.8pt, arc=2mm]

    \textbf{Dimension concernée :} gold\_dim\_staff.

    Un Slowly Changing Dimension (SCD) permet de gérer les évolutions
    des attributs descriptifs d’une dimension dans le temps.

    \textbf{Deux stratégies principales :}

    \textbf{SCD Type 1 :}
    \begin{itemize}
    \item Les valeurs sont écrasées lors d’un changement.
    \item Pas d’historique conservé.
    \item Modèle simple et peu volumineux.
    \item Adapté si l’historique n’a pas d’intérêt analytique.
    \end{itemize}

    \textbf{SCD Type 2 :}
    \begin{itemize}
    \item Création d’une nouvelle ligne lors d’un changement.
    \item Ajout de colonnes de validité (valid\_from, valid\_to, is\_current).
    \item Conservation complète de l’historique.
    \item Permet des analyses temporelles fiables.
    \end{itemize}

    \textbf{Choix recommandé : SCD Type 2}

    Justification :
    \begin{itemize}
    \item Un membre du personnel peut changer de rôle ou de service.
    \item Le directeur peut vouloir analyser l’impact organisationnel :
    \begin{itemize}
    \item Le taux de refus a-t-il augmenté après un changement d’équipe ?
    \item La satisfaction a-t-elle évolué après une réorganisation ?
    \end{itemize}
    \item Sans historisation, ces analyses seraient biaisées.
    \end{itemize}

    Dans le cadre pédagogique du TP2 (volumétrie faible et complexité limitée),
    un SCD Type 1 peut être toléré pour simplifier l’implémentation.
    Cependant, d’un point de vue analytique et architectural,
    le SCD Type 2 est le choix le plus robuste.

    \end{tcolorbox}

    \begin{tcolorbox}[colback=blue!5, colframe=blue!60, title=\textbf{Commentaire}]
    Le choix d’un SCD Type 1 en SILVER simplifie le modèle,
    mais limite les capacités analytiques.

    Dans un contexte réel, la décision d’historiser (Type 2)
    dépend du besoin métier et du coût de complexité supplémentaire.
    \end{tcolorbox}
\end{enumerate}


% ---------------------------------------------------------------------------------------------

\section{Étape 6 : Exploitation décisionnelle (Dashboard BI)}

\subsection*{Objectif}

Valider la cohérence du datamart GOLD en construisant un premier
tableau de bord décisionnel permettant de répondre aux questions
du directeur.

\subsection*{Outil recommandé}

Power BI (ou outil équivalent : Tableau, Superset, etc.).

\subsection*{Travail demandé}

\begin{enumerate}
  \item Se connecter à la base \file{clinic\_dwh.db}.
  \item Importer uniquement les tables \file{gold\_*}.
  \item Vérifier les relations entre tables (PK/FK).
\end{enumerate}

\subsection*{KPI minimums attendus}

Construire au minimum les indicateurs suivants :

\begin{itemize}
  \item Taux d’occupation moyen par service.
  \item Nombre de patients refusés par semaine.
  \item Taux de refus.
  \item Satisfaction moyenne.
  \item Durée moyenne de séjour.
  \item Nombre de membres du personnel présents.
\end{itemize}

Pour chaque KPI, préciser :

\begin{itemize}
  \item la définition métier,
  \item la formule utilisée,
  \item le niveau d’agrégation.
\end{itemize}

\subsection*{Structure attendue du dashboard}

Créer au minimum :

\begin{itemize}
  \item Une page \textbf{Pilotage global} :
  \begin{itemize}
    \item KPI principaux (cartes),
    \item évolution temporelle,
    \item filtre par service.
  \end{itemize}
  \item Une page \textbf{Analyse détaillée} :
  \begin{itemize}
    \item comparaison entre services,
    \item focus sur les refus et la satisfaction,
    \item possibilité de filtrer par période.
  \end{itemize}
\end{itemize}

\subsection*{Analyse attendue}

Rédiger une courte analyse (10 lignes environ) répondant à au moins
\textbf{3 questions du directeur} à partir du dashboard construit.

\subsection*{Livrables attendus}

\begin{itemize}
  \item Fichier Power BI (\file{.pbix}) ou captures d’écran.
  \item Liste des KPI définis (avec formules).
  \item Analyse rédigée.
\end{itemize}

\begin{tcolorbox}[title=\textbf{Validation décisionnelle via Power BI :}, colback=gray!5, colframe=black, boxrule=0.8pt, arc=2mm]

\textbf{Objectif :}
Vérifier que le datamart permet réellement de répondre
aux questions stratégiques.

\textbf{Bonnes pratiques :}
\begin{itemize}
\item Importer uniquement les tables gold\_*.
\item Vérifier le schéma en étoile.
\item Créer des mesures DAX (pas utiliser directement les colonnes).
\item Centraliser les mesures dans une table dédiée.
\end{itemize}

\textbf{Exemples d’analyses attendues :}
\begin{itemize}
\item Identification des services en tension.
\item Corrélation entre staff et taux de refus.
\item Évolution de la satisfaction en période de forte activité.
\end{itemize}

Si le dashboard répond clairement à au moins 3 questions du directeur,
le modèle GOLD est considéré comme valide.

\end{tcolorbox}

\end{document}